\SourceLeaf{403}%
\[
\left.\begin{aligned}
\frac{r-q}{p^*-q}&=\frac{(a-b)(d-c)(ab-cd)^2}{(a-d)(b-c)(ad-bc)^2}.pp'\\
\frac{r-q}{p^*-r}&=\frac{(a-b)(d-c)(ab-cd)^2}{(a-c)(b-d)(ac-bd)^2}.pp'\\
\frac{p^*-r}{p^*-q}&=\frac{(a-c)(b-d)(ac-bd)^2}{(a-d)(b-c)(ad-bc)^2}
\end{aligned}\right\}\ \text{are real numbers}
\]
or \(p^*,q,r\) lie on a straight line perpendicular to \(0p^*p\)
\par\medskip\Needspace{5\baselineskip}
\pdfbookmark[2]{V.}{carnot-zusatz-v}
\centerline{\textbf{V.}}\nobreak\smallskip
\{Inscribe in a given circle a polygon whose sides pass through an equal number of given points. \textsc{Schumacher}.\}

Let \(r\) be the radius of the circle and let the coordinates of the vertices of the polygon be
\[
\begin{array}{llll}
r\cos\varphi,&r\cos\varphi',&r\cos\varphi''&\text{etc.}\\
r\sin\varphi,&r\sin\varphi',&r\sin\varphi''&\text{etc.}
\end{array}
\]
Finally, let the coordinates of the given points through which the sides of the polygon pass when extended (the sides joining the first and second vertices, the second and third, and so on, respectively) be
\[
\begin{array}{llll}
a\cos A,&a'\cos A',&a''\cos A''&\text{etc.}\\
a\sin A,&a'\sin A',&a''\sin A''&\text{etc.}
\end{array}
\]
\par
When three points whose coordinates are \(x,y;\ x',y';\ x'',y''\) lie on a straight line, the condition
\[
xy'+x'y''+x''y-x'y-x''y'-xy''=0
\]
holds. It follows that
\[
\left.\begin{aligned}
rr\cos\varphi\sin\varphi'+ra\cos\varphi'\sin A+ar\cos A\sin\varphi\\
-rr\cos\varphi'\sin\varphi-ar\cos A\sin\varphi'-ra\cos\varphi\sin A
\end{aligned}\right\}=0
\]
or
\[
r\sin(\varphi'-\varphi)-a\sin(\varphi'-A)+a\sin(\varphi-A)=0
\]
or
\[
r\cos\tfrac12(\varphi'-\varphi)=a\cos\tfrac12(\varphi'+\varphi-2A)
\]
